\chapter{Cointegração e VECM}
\label{cap:coint}
% ── índice ──────────────────────────────────────────────────────
\index{cointegração|textbf}
\index{VECM|textbf}
\index{modelo de correção de erros|see{VECM}}
\index{Johansen, teste de|textbf}
\index{johansen@\texttt{johansen()}|textbf}
\index{vecm@\texttt{vecm()}|textbf}
\index{relação de longo prazo}

% ─────────────────────────────────────────────────────────────────────────────
\section{Relação de longo prazo entre séries I(1)}

Duas séries $I(1)$ são \textbf{cointegradas} se existe uma combinação linear
entre elas que é $I(0)$. Formalmente, $y_t$ e $x_t$ são cointegradas se
existe $\beta$ tal que:
\[
  u_t = y_t - \beta x_t \sim I(0)
\]

A intuição é a de uma \emph{correia}: séries que individualmente seguem
passeios aleatórios podem ser atraídas por uma força de equilíbrio de longo
prazo. Preços de ativos no mesmo mercado, taxa de câmbio e paridade de poder
de compra, juros de curto e longo prazo são exemplos clássicos.

Séries $I(1)$ não cointegradas divergem indefinidamente — a regressão em
nível produz \textbf{regressão espúria}: $R^2$ alto e $t$-estatísticas
significativas sem relação causal real.

% ─────────────────────────────────────────────────────────────────────────────
\section{Teste de Engle-Granger}

O procedimento de Engle-Granger (1987) testa cointegração em dois passos:

\begin{enumerate}
  \item Estima a \textbf{relação de longo prazo} por OLS:
        $\hat u_t = y_t - \hat\beta x_t$
  \item Aplica o \textbf{teste ADF} nos resíduos $\hat u_t$ — sob $H_0$
        os resíduos têm raiz unitária (sem cointegração).
\end{enumerate}

Os valores críticos do teste são mais negativos do que os do ADF padrão,
pois $\hat u_t$ é resíduo estimado, não série diretamente observada.

% ─────────────────────────────────────────────────────────────────────────────
\section{Modelo de Correção de Erros (VECM)}

Se $y_t$ e $x_t$ são cointegradas, o VECM (\emph{Vector Error Correction
Model}) decompõe a dinâmica em componente de curto prazo e ajuste ao
equilíbrio de longo prazo:

\[
  \Delta \mathbf{y}_t
  = \boldsymbol{\alpha} \underbrace{(\boldsymbol{\beta}' \mathbf{y}_{t-1})}_{\text{desvio do equilíbrio}}
  + \sum_{j=1}^{p-1} \mathbf{\Gamma}_j \Delta \mathbf{y}_{t-j}
  + \boldsymbol{\varepsilon}_t
\]

\begin{itemize}
  \item $\boldsymbol{\beta}$ — \textbf{vetor cointegrante}: define a relação
        de equilíbrio de longo prazo $\boldsymbol{\beta}'\mathbf{y}_t = 0$.
  \item $\boldsymbol{\alpha}$ — \textbf{coeficientes de ajuste}: velocidade
        com que cada variável corrige desvios do equilíbrio. Sinal negativo
        em $\alpha_y$ significa que $y$ cai quando está acima do equilíbrio.
  \item $\mathbf{\Gamma}_j$ — dinâmica de curto prazo.
\end{itemize}

O VECM é estimado por ML via procedimento de Johansen, que testa o número
de vetores cointegrantes (\emph{rank}) via estatística de traço (\emph{trace})
ou máximo autovalor.

% ─────────────────────────────────────────────────────────────────────────────
\section{Verificação por DGP}

\subsection*{Caso 1 — cointegrado: $y_t = 2 x_t + \varepsilon_t$}

$x_t$ é um passeio aleatório ($I(1)$); $\varepsilon_t \sim I(0)$. Logo
$y_t - 2 x_t = \varepsilon_t \sim I(0)$: as séries são cointegradas com
vetor $[1,\; -2]$.

\begin{lstlisting}[caption={DGP cointegrado e teste de Engle-Granger}]
seed(42)
input df
id
1
end
for i in 1..=9 { df = append(df, df) }
df |> filter(_n <= 300)
generate df e1 = rnormal()
generate df e2 = rnormal()
generate df x  = cumsum(e1)
generate df y  = 2.0*x + e2
coint(df, y, x)
\end{lstlisting}

\begin{lstlisting}[language={}, basicstyle=\ttfamily\small,
  frame=none, numbers=none, backgroundcolor=\color{codbg},
  xleftmargin=0pt, xrightmargin=0pt, breaklines=false]
============= Engle-Granger Cointegration Test =============
  Variables: y, x
  H₀: sem cointegração
  ADF statistic:    -6.4808
  Critical values:  1%=-3.900  5%=-3.340  10%=-3.040
  Vetor cointegrante: [0.4778, 2.0002]
  Conclusion: REJEITA H₀ — séries cointegradas
\end{lstlisting}

A estatística $-6{,}48$ está muito abaixo do valor crítico de 1\%
($-3{,}90$). O vetor cointegrante estimado $[0{,}48;\; 2{,}00]$ recupera
$\beta = 2{,}0$ com precisão.

\subsection*{Caso 2 — não cointegrado: dois passeios aleatórios independentes}

\begin{lstlisting}[caption={Dois random walks — sem cointegração}]
generate df y = cumsum(e2)
coint(df, y, x)
\end{lstlisting}

\begin{lstlisting}[language={}, basicstyle=\ttfamily\small,
  frame=none, numbers=none, backgroundcolor=\color{codbg},
  xleftmargin=0pt, xrightmargin=0pt, breaklines=false]
============= Engle-Granger Cointegration Test =============
  Variables: y, x
  H₀: sem cointegração
  ADF statistic:    -3.1058
  Critical values:  1%=-3.900  5%=-3.340  10%=-3.040
  Conclusion: Não rejeita H₀ — sem cointegração
\end{lstlisting}

\subsection*{VECM — dinâmica de ajuste}

Com cointegração confirmada, estima-se o VECM:

\begin{lstlisting}[caption={VECM com rank 1}]
generate df y = 2.0*x + e2
let m = vecm(df, y, x, lags=1)
display m
\end{lstlisting}

\begin{lstlisting}[language={}, basicstyle=\ttfamily\small,
  frame=none, numbers=none, backgroundcolor=\color{codbg},
  xleftmargin=0pt, xrightmargin=0pt, breaklines=false]
======================== VECM (Johansen ML) - Rank 1 =========================
No. Variables:                2
Observations:               299

------------------------ Cointegration Vectors (Beta) ------------------------
[ 3.6175, -7.2357 ]  →  normalizado: [ 1, -2.000 ]

---------------------- Adjustment Coefficients (Alpha) -----------------------
  y:  -0.3133   (corrige desvios do equilíbrio)
  x:  -0.0179   (quase exógena — ajuste mínimo)

----------------------- Johansen Eigenvalues (Lambda) ------------------------
    λ₁ = 0.5017     λ₂ = 0.0004
==============================================================================
\end{lstlisting}

O vetor cointegrante normalizado $[1,\; -2{,}00]$ recupera o verdadeiro
$\beta = 2{,}0$. O coeficiente de ajuste $\hat\alpha_y = -0{,}313$: quando
$y_{t-1} > 2 x_{t-1}$ (acima do equilíbrio), $\Delta y_t$ é negativo —
$y$ se corrige para baixo. $x$ quase não ajusta ($\hat\alpha_x \approx -0{,}02$),
pois é a variável exógena que define o equilíbrio.

% ─────────────────────────────────────────────────────────────────────────────
\section{Sintaxe}

\begin{lstlisting}
coint(df, y, x)
coint(df, y, x, lags=4)
vecm(df, y, x, lags=1)
vecm(df, y1, y2, y3, lags=2)
\end{lstlisting}

\begin{lstlisting}[caption={Fluxo típico — taxa de câmbio e paridade}]
load "cambio.csv" as df

adf(df, spot)
adf(df, forward)

coint(df, spot, forward)

// se cointegradas:
let m = vecm(df, spot, forward, lags=2)
display m
\end{lstlisting}

\begin{nota}
  O VECM pressupõe que todas as variáveis são $I(1)$ e cointegradas.
  Verifique com ADF antes de estimar. Se as séries forem $I(1)$ mas não
  cointegradas, use o VAR em primeiras diferenças (cap.~\ref{cap:var}).
\end{nota}

\section{Validação empírica}

O DGP cointegrado deste capítulo é usado no caso de validação
\texttt{coint_book} em \texttt{validation/cases/}. O caso estima um VECM de
posto 1 em \hayashi{} e o compara com implementações de referência em R e
Python. Execute \lstinline{hay validate} para ver o estado atual.
